\chapter{实解析函数}
\label{chap:real-analytic-functions-solutions}

\begin{solution}{analytic-coefficients-derivatives}
逐项求导并在 \(x=a\) 处代入，得到
\[
 c_0=f(a),\qquad c_1=f'(a),\qquad
 c_2=\frac{f''(a)}{2!},\qquad
 c_3=\frac{f'''(a)}{3!}.
\]
一般地 \(c_n=f^{(n)}(a)/n!\)。
\end{solution}

\begin{solution}{analytic-recentering-geometric}
写
\[
 \frac1{1-x}
 =\frac1{1-c-(x-c)}
 =\frac1{1-c}\frac1{1-\dfrac{x-c}{1-c}}.
\]
因此
\[
 \frac1{1-x}
 =\sum_{n=0}^\infty\frac{(x-c)^n}{(1-c)^{n+1}},
 \qquad |x-c|<1-c.
\]
新中心为 \(c\)，收敛半径为 \(1-c\)，它等于中心到实奇点 \(1\) 的距离。
\end{solution}

\begin{solution}{analytic-recentering-absolute-regrouping}
对 \(x=c+h\)，二项展开后的项为
\[
 a_n\binom nk(c-x_0)^{n-k}h^k,\qquad0\le k\le n.
\]
其绝对值的全部总和是
\[
\begin{aligned}
 \sum_{n=0}^\infty\sum_{k=0}^n
 |a_n|\binom nk|c-x_0|^{n-k}|h|^k
 &=\sum_{n=0}^\infty|a_n|
 (|c-x_0|+|h|)^n.
\end{aligned}
\]
当 \(|c-x_0|+|h|<R\) 时最后的幂级数收敛。绝对收敛的二重级数可按
\(k\) 重新求和，所以交换次序和按 \(h^k\) 合并都合法。
\end{solution}

\begin{solution}{analytic-taylor-remainder-equivalence}
若在 \(|x-a|<r\) 内对每个 \(x\) 都有 \(R_{N,a}f(x)\to0\)，则
\[
 f(x)=\lim_{N\to\infty}T_{N,a}f(x)
 =\sum_{n=0}^\infty\frac{f^{(n)}(a)}{n!}(x-a)^n.
\]
这就是 \(a\) 处的解析表示。反过来，一旦得到幂级数表示，幂级数在每个
\(|x-a|\le\rho<r\) 上一致收敛，故 Taylor 余项在这些内闭区间上一致
趋零。
\end{solution}

\begin{solution}{analytic-derivative-bound-exponential}
固定 \(a\) 和 \(r>0\)。在 \((a-r,a+r)\) 上，
\[
 |(e^x)^{(n)}|=e^x\le e^{a+r}\le e^{a+r}\,1^n n!.
\]
所以可取 \(M=e^{a+r}\)、\(A=1\)。导数增长定理先在
\(|x-a|<\min\{r,1\}\) 内给出 Taylor 表示。这个固定 \(A=1\) 的估计
只保证半径至少1；要得到整个实轴上的等式，直接使用余项
\[
 |R_N(x)|\le e^{\max\{a,x\}}\frac{|x-a|^{N+1}}{(N+1)!}\to0,
\]
可知 \(e^x\) 在任意 \(x\in\R\) 处等于其 Taylor 级数。
\end{solution}

\begin{solution}{smooth-derivatives-bounded-not-enough}
闭区间上每个固定的连续函数 \(f^{(n)}\) 都有界，但所得上界可随 \(n\)
任意快地增长。Lagrange 余项需要比较
\[
 \frac{\sup|f^{(N+1)}|}{(N+1)!}|x-a|^{N+1}
\]
随 \(N\) 的行为；逐阶存在有限上界并不能使该量趋零。平坦函数
\(\phi(x)=e^{-1/x^2}\)（\(x\ne0\)）、\(\phi(0)=0\) 在任意闭区间上
每阶导数都连续有界，却在零点所有 Taylor 系数为零而函数不恒为零，故不
解析。
\end{solution}

\begin{solution}{analytic-product-coefficients}
在共同收敛邻域内两个幂级数绝对收敛，因此二重级数
\[
 \sum_{j,k\ge0}a_jb_kh^{j+k}
\]
绝对收敛，可以按对角线 \(j+k=n\) 分组。于是
\[
 f(h)g(h)=\sum_{n=0}^\infty
 \left(\sum_{k=0}^na_kb_{n-k}\right)h^n.
\]
绝对收敛正用于保证重排和分组不改变总和。
\end{solution}

\begin{solution}{analytic-reciprocal-local}
令
\[
 u(x)=\frac{f(x)-f(a)}{f(a)}.
\]
它在 \(a\) 处解析且 \(u(a)=0\)。缩小邻域可使 \(|u(x)|<q<1\)，从而
\[
 \frac1{f(x)}
 =\frac1{f(a)}\frac1{1+u(x)}
 =\frac1{f(a)}\sum_{m=0}^\infty(-u(x))^m.
\]
若 \(u(a+h)=\sum_{n\ge1}c_nh^n\)，再缩小 \(r\) 使
\(\sum_{n\ge1}|c_n|r^n<1\)。把各次幂展开后，所有项的绝对值由一个
几何级数控制，故可按 \(h\) 的次数分组为收敛幂级数。这直接证明倒数
解析。
\end{solution}

\begin{solution}{analytic-quotient}
对任意 \(a\in I\)，因 \(g(a)\ne0\)，解析倒数定理说明 \(1/g\) 在
\(a\) 的某邻域解析。解析函数的乘积仍解析，所以
\[
 \frac fg=f\cdot\frac1g
\]
在 \(a\) 附近解析。因 \(a\) 任意，商在整个 \(I\) 上解析。
\end{solution}

\begin{solution}{analytic-composition-majorant}
写
\[
 u(h)=\sum_{n\ge1}a_nh^n,\qquad
 g(f(a)+y)=\sum_{m\ge0}b_my^m.
\]
将 \(u(h)^m\) 展开后得到
\[
 \sum_{m=0}^\infty b_m
 \sum_{n_1,\ldots,n_m\ge1}
 a_{n_1}\cdots a_{n_m}h^{n_1+\cdots+n_m}.
\]
取 \(r\) 使 \(A(r)=\sum_{n\ge1}|a_n|r^n\) 落在外层幂级数的绝对
收敛区间。上述三级求和在 \(|h|\le r\) 上的绝对值总和不超过
\[
 \sum_{m=0}^\infty|b_m|A(r)^m<\infty.
\]
故可按总次数 \(n_1+\cdots+n_m\) 分组，得到关于 \(h\) 的绝对收敛
幂级数。
\end{solution}

\begin{solution}{analytic-log-series}
在 \(|x|<1\) 内，
\[
 \frac1{1+x}=\sum_{n=0}^\infty(-1)^nx^n.
\]
在从 \(0\) 到 \(x\) 的内闭区间上逐项积分，
\[
 \log(1+x)
 =\int_0^x\frac{dt}{1+t}
 =\sum_{n=0}^\infty\frac{(-1)^nx^{n+1}}{n+1}.
\]
重新编号即得题中公式。
\end{solution}

\begin{solution}{analytic-rational-expansion}
由几何级数，
\[
 \frac1{1+x^2}
 =\sum_{n=0}^\infty(-1)^nx^{2n}
 \qquad(|x|<1).
\]
系数根值上极限为 \(1\)，故收敛半径为 \(1\)。函数在整个实轴分母不为
零，因而 \(C^\infty\)。但它在零点的 Taylor 系数已经由唯一性确定为
上述系数，而该级数在 \(|x|>1\) 时通项不趋零，所以不能用这一单一展开
表示整个实轴。
\end{solution}

\begin{solution}{analytic-zero-order-product}
写
\[
 f(x)=(x-a)^mF(x),\qquad
 g(x)=(x-a)^nG(x),
\]
其中 \(F,G\) 解析且 \(F(a)G(a)\ne0\)。则
\[
 f(x)g(x)=(x-a)^{m+n}F(x)G(x),
\]
而 \(FG\) 在 \(a\) 处非零，故乘积零点阶为 \(m+n\)。
\end{solution}

\begin{solution}{analytic-zero-order-derivative}
写 \(f(x)=(x-a)^mg(x)\)，其中 \(g(a)\ne0\)。则
\[
 f'(x)=(x-a)^{m-1}
 \bigl(mg(x)+(x-a)g'(x)\bigr).
\]
括号内函数在 \(a\) 处取值 \(mg(a)\ne0\)，所以 \(f'\) 的消失阶恰为
\(m-1\)。当 \(m=1\) 时括号内在 \(a\) 处非零，\(f'(a)\ne0\)，
称消失阶为0，不能称 \(a\) 为 \(f'\) 的零点。
\end{solution}

\begin{solution}{analytic-zeros-compact-finite}
若 \([a,b]\) 中有无穷多个零点，紧致性和 Bolzano--Weierstrass 定理
给出一个收敛零点子列，其极限 \(c\in[a,b]\)。由于函数在包含
\([a,b]\) 的开区间上解析，\(c\) 是定义域内部的零点聚点。恒等原理迫使
\(f\) 在整个连通开区间上恒为零，与假设矛盾。因此零点只能有有限多个。
\end{solution}

\begin{solution}{analytic-identity-data-set}
差函数 \(h=f-g\) 在 \(I\) 上解析，并满足 \(h(x_n)=0\)。
\(x_n\to a\in I\)、\(x_n\ne a\) 说明其零点在定义域内部有聚点。由
恒等原理 \(h\equiv0\)，即 \(f\equiv g\)。
\end{solution}

\begin{solution}{flat-function-derivative-form}
初始 \(P_0(t)=1\)。若
\[
 \phi^{(n)}(x)=P_n(1/x)e^{-1/x^2},
\]
则令 \(t=1/x\)，有
\[
\begin{aligned}
 \phi^{(n+1)}(x)
 &=\left[-x^{-2}P_n'(1/x)+2x^{-3}P_n(1/x)\right]e^{-1/x^2}.
\end{aligned}
\]
所以
\[
 P_{n+1}(t)=-t^2P_n'(t)+2t^3P_n(t),
\]
仍是多项式。归纳完成。
\end{solution}

\begin{solution}{smooth-bump-not-analytic}
在 \(|x|<1\) 内有 \(1-x^2>0\)，函数 \(x\mapsto1/(1-x^2)\) 光滑；
与指数函数复合后，\(\psi\) 在该区间内光滑。靠近 \(x=1\) 时，
\[
 1-x^2=(1-x)(1+x)
\]
与 \(1-x\) 同阶；靠近 \(x=-1\) 时与 \(1+x\) 同阶。反复求导得到的
每一阶导数都是 \(Q_n(x,t)e^{-t}\)，其中 \(t=(1-x^2)^{-1}\)，
\(Q_n\) 为二元多项式。由链式法则，初值和递推为
\[
 Q_0=1,\qquad
 Q_{n+1}=\partial_xQ_n+2xt^2\partial_tQ_n-2xt^2Q_n.
\]
每个 \(Q_n\) 在 \(|x|<1\) 时的绝对值不超过某个 \(t\) 的多项式，
因此 \(Q_n(x,t)e^{-t}\to0\)。在 \(a=\pm1\) 处还需检查差商：
内侧 \(|x-a|^{-1}=(1+|x|)t\le2t\)，所以
\[
 \left|\frac{Q_n(x,t)e^{-t}}{x-a}\right|
 \le2t|Q_n(x,t)|e^{-t}\longrightarrow0.
\]
外侧为零。归纳得到端点所有阶导数存在、等于零并连续，故
\(\psi\in C^\infty(\R)\)。

在 \(x=1\) 处所有 Taylor 系数为零，但左侧邻域函数为正，故不解析；
\(x=-1\) 同理。也可反设它在边界解析：它在边界外侧一个开区间恒为零，
恒等原理便迫使它在包含边界的连通邻域恒为零，与内侧正值矛盾。
\end{solution}

\begin{solution}{analytic-compact-support}
设 \(f\) 在连通开区间 \(I\) 上解析，且其支集 \(K\) 是 \(I\) 的紧子集。
非空开区间不可能为紧集，故 \(K\ne I\)。紧集闭，因而 \(I\setminus K\)
是非空开集，含有一个开子区间；\(f\) 在该子区间恒为零。其零点有内部
聚点，恒等原理给出 \(f\equiv0\)。

\(C^\infty\) 范畴没有恒等原理。正文的 \(\psi\) 是非零、光滑且支集
包含在 \([-1,1]\) 中的函数，直接给出反例。
\end{solution}
